

\section{חיבור תנע זוויתי - חזרה והעמקה}

\begin{mainbox}{מוטיבציה פיזיקלית}
במערכות רב-חלקיקיות (למשל, אטום עם מספר אלקטרונים, או צימוד מסילה-ספין), הגודל הפיזיקלי השמור אינו בהכרח התנע הזוויתי של כל חלקיק בנפרד ($\mathbf{J}_1, \mathbf{J}_2$), אלא התנע הזוויתי הכולל של המערכת $\mathbf{J} = \mathbf{J}_1 + \mathbf{J}_2$. לכן, נרצה לעבור מתיאור המערכת באמצעות התנעים הנפרדים לתיאור באמצעות התנע הכולל.
\end{mainbox}

\subsection{מרחב המכפלה ומרחב התנע הכולל}

נתונות שתי מערכות עם תנע זוויתי $j_1$ ו-$j_2$. המרחב ההילברט הכולל הוא מכפלה טנזורית של המרחבים:
\[
V = V_{j_1} \otimes V_{j_2}
\]
הממד של המרחב הוא $(2j_1+1)(2j_2+1)$. קיימים שני בסיסים טבעיים לתיאור המערכת:

\begin{enumerate}
    \item \textbf{הבסיס הלא-מצומד (\LR{Product Basis}):}
    בסיס זה נפרש על ידי המצבים העצמיים של האופרטורים $\{J_1^2, J_{1z}, J_2^2, J_{2z}\}$. המצבים מסומנים כ:
    \[
    |j_1 m_1; j_2 m_2\rangle \equiv |j_1 m_1\rangle \otimes |j_2 m_2\rangle
    \]
    
    \item \textbf{הבסיס המצומד (\LR{Coupled Basis}):}
    בסיס זה נפרש על ידי המצבים העצמיים של האופרטורים $\{J^2, J_z, J_1^2, J_2^2\}$, כאשר $\mathbf{J} = \mathbf{J}_1 + \mathbf{J}_2$. המצבים מסומנים כ:
    \[
    |J M; j_1 j_2\rangle
    \]
\end{enumerate}

\begin{theorembox}{פירוק להצגות בלתי-פריקות}
המעבר בין הבסיסים שקול לפירוק מרחב המכפלה לסכום ישר של תתי-מרחבים אינווריאנטיים תחת סיבוב (הצגות בלתי פריקות של חבורת הסיבובים):
\[
\mathbf{j}_1 \otimes \mathbf{j}_2 = \bigoplus_{J=|j_1-j_2|}^{j_1+j_2} \mathbf{J}
\]
כל ערך של $J$ מופיע בדיוק פעם אחת בפירוק זה.
\end{theorembox}

\subsection{מקדמי קלבש-גורדון (\LR{Clebsch-Gordan Coefficients})}

הקשר בין הבסיסים הוא טרנספורמציה אוניטרית (שינוי בסיס). המקדמים בטרנספורמציה זו נקראים מקדמי קלבש-גורדון:

\[
|J M\rangle = \sum_{m_1, m_2} \langle j_1 m_1; j_2 m_2 | J M \rangle |j_1 m_1; j_2 m_2\rangle
\]

\begin{notebox}{תנאי סלקציה למקדמים}
המקדמים $\langle j_1 m_1; j_2 m_2 | J M \rangle$ שונים מאפס רק אם מתקיימים התנאים הבאים:
\begin{itemize}
    \item שימור היטל התנע הזוויתי: $M = m_1 + m_2$.
    \item אי-שוויון המשולש: $|j_1 - j_2| \le J \le j_1 + j_2$.
\end{itemize}
\end{notebox}

\subsubsection{אלגוריתם לחישוב המקדמים}

תהליך מציאת המקדמים מתבצע באופן שיטתי, החל מהמצב בעל המשקל הגבוה ביותר (\LR{Highest Weight State}):

\begin{enumerate}
    \item \textbf{זיהוי המצב המקסימלי:} המצב עם $M_{max} = j_1 + j_2$ הוא יחיד. הוא חייב להיות המצב העליון של ההצגה עם $J_{max} = j_1 + j_2$. לכן:
    \[
    |J_{max}, M_{max}\rangle = |j_1 j_1; j_2 j_2\rangle
    \]
    (המקדם הוא 1, לפי קונבנציית פאזה של \LR{Condon-Shortley}).
    
    \item \textbf{הפעלת אופרטורי הורדה:} מפעילים את אופרטור ההורדה הכולל $J_- = J_{1-} + J_{2-}$ על שני אגפי המשוואה.
    \[
    J_- |J M\rangle = \hbar\sqrt{J(J+1) - M(M-1)} |J, M-1\rangle
    \]
    כך מקבלים את כל המצבים בהצגה $J_{max}$.
    
    \item \textbf{אורתוגונליות:} כדי למצוא את המצב העליון של ההצגה הבאה ($J = J_{max}-1$), מחפשים קומבינציה לינארית של מצבי הבסיס הלא-מצומד עם $M = J_{max}-1$ שתהיה אורתוגונלית למצב $|J_{max}, J_{max}-1\rangle$ שמצאנו בשלב הקודם.
    
    \item חוזרים על התהליך עבור כל ערכי $J$ האפשריים.
\end{enumerate}

\begin{center}
\begin{tikzpicture}[scale=0.8]
    \begin{axis}[
        title={מרחב המצבים ($m_1, m_2$)},
        xlabel={$m_1$},
        ylabel={$m_2$},
yo        xmin=-2.5, xmax=2.5,
        ymin=-2.5, ymax=2.5,
        grid=both,
        axis lines=middle,
        xtick={-2,-1,0,1,2},
        ytick={-2,-1,0,1,2}
    ]
        % Example points for j1=2, j2=1 (Conceptual representation)
        \addplot[only marks, mark=*, color=blue] coordinates {
            (2,1) (1,1) (0,1) (-1,1) (-2,1)
            (2,0) (1,0) (0,0) (-1,0) (-2,0)
            (2,-1) (1,-1) (0,-1) (-1,-1) (-2,-1)
        };
        \node[anchor=south west] at (axis cs:2,1) {המשקל הגבוה ביותר ($J=j_1+j_2$)};
    \end{axis}
\end{tikzpicture}
\end{center}

\newpage

\section{אופרטורים טנזוריים (\LR{Tensor Operators})}

עד כה עסקנו בטרנספורמציה של מצבים (וקטורים במרחב הילברט) תחת סיבוב. כעת נדון בטרנספורמציה של \textbf{אופרטורים}.

\subsection{אופרטור סקלרי ווקטורי}

\begin{definitionbox}{אופרטור סקלרי}
אופרטור $S$ נקרא \textbf{סקלרי} אם הוא אינו משתנה תחת סיבוב המערכת. פורמלית, אם $U(R)$ הוא אופרטור הסיבוב במרחב הילברט המתאים לסיבוב מרחבי $R$:
\[
S' = U(R) S U^\dagger(R) = S
\]
דוגמאות: $H$ (המילטוניאן עם סימטריה כדורית), $\mathbf{r}^2$, $\mathbf{p}^2$.
\end{definitionbox}

\begin{definitionbox}{אופרטור וקטורי}
אוסף של שלושה אופרטורים $\mathbf{V} = (V_x, V_y, V_z)$ נקרא \textbf{אופרטור וקטורי} אם הרכיבים עוברים טרנספורמציה כמו וקטור קרטזי רגיל תחת סיבוב:
\[
U(R) V_i U^\dagger(R) = \sum_{j} R_{ij} V_j
\]
\end{definitionbox}
\subsection{טנזורים ספריים (\LR{Spherical Tensors})}

ההגדרה הקרטזית של אופרטור וקטורי אינה נוחה לשימוש בבעיות קוונטיות עם סימטריה כדורית (כמו חישוב אלמנטי מטריצה בין מצבי תנע זוויתי). לשם כך נגדיר טנזורים ספריים.

\subsubsection{הקשר לבסיס הספירי}
אנו יודעים כי ההרמוניות הספריות $Y_{l}^m$ מהוות בסיס להצגה הבלתי פריקה מסדר $l$ של חבורת הסיבובים. בפרט, עבור $l=1$, יש לנו שלושה מצבים המקבילים לווקטור.
ניתן לעבור מהבסיס הקרטזי $(x, y, z)$ לבסיס הספירי $(r_{+1}, r_0, r_{-1})$ באמצעות הטרנספורמציה:

\begin{resultbox}{מעבר בסיס: קרטזי לספירי}
\[
r_{+1} = -\frac{x + iy}{\sqrt{2}}, \quad r_0 = z, \quad r_{-1} = \frac{x - iy}{\sqrt{2}}
\]
\end{resultbox}
באופן דומה מגדירים אופרטורים טנזוריים ספריים מדרגה $k=1$:
\[
V_{+1}^{(1)} = -\frac{V_x + iV_y}{\sqrt{2}}, \quad V_0^{(1)} = V_z, \quad V_{-1}^{(1)} = \frac{V_x - iV_y}{\sqrt{2}}
\]

\subsubsection{הגדרה כללית}
טנזור ספירי מסדר $k$, המסומן $T^{(k)}$, הוא סט של $2k+1$ אופרטורים $T_q^{(k)}$ (כאשר $q = -k, \dots, k$) אשר עוברים טרנספורמציה תחת סיבוב לפי ההצגה הבלתי פריקה מסדר $k$:

\[
U(R) T_q^{(k)} U^\dagger(R) = \sum_{q'=-k}^{k} D_{q'q}^{(k)}(R) T_{q'}^{(k)}
\]
כאשר $D_{q'q}^{(k)}$ הן מטריצות ויגנר (\LR{Wigner D-matrices}).

\subsection{פירוק מכפלה של אופרטורים וקטוריים}

כאשר מכפילים שני אופרטורים וקטוריים $X_i$ ו-$Y_j$, מקבלים טנזור קרטזי מדרגה שנייה $T_{ij} = X_i Y_j$ בעל 9 רכיבים. כמו בחיבור תנע זוויתי $1 \otimes 1 = 0 \oplus 1 \oplus 2$, ניתן לפרק טנזור זה לרכיבים אינווריאנטיים (בלתי פריקים):

\begin{equation}
X_i Y_j = \underbrace{\frac{1}{3}(\mathbf{X}\cdot\mathbf{Y})\delta_{ij}}_{\text{סקלר } (k=0)} + \underbrace{\frac{1}{2}(X_i Y_j - X_j Y_i)}_{\text{וקטור/אנטי-סימטרי } (k=1)} + \underbrace{\left(\frac{X_i Y_j + X_j Y_i}{2} - \frac{1}{3}(\mathbf{X}\cdot\mathbf{Y})\delta_{ij}\right)}_{\text{סימטרי חסר עקבה } (k=2)}
\end{equation}

\begin{itemize}
    \item \textbf{החלק הסקלרי ($k=0$):} קשור לטרייס (עקבה) של הטנזור, $\mathbf{X}\cdot\mathbf{Y}$. מכיל דרגת חופש אחת.
    \item \textbf{החלק הוקטורי ($k=1$):} החלק האנטי-סימטרי, שקול למכפלה וקטורית $\mathbf{X} \times \mathbf{Y}$. מכיל 3 דרגות חופש.
    \item \textbf{החלק הטנזורי ($k=2$):} החלק הסימטרי וחסר העקבה (\LR{Traceless Symmetric}). זהו למשל אופרטור הקוודרופול. מכיל 5 דרגות חופש.
\end{itemize}

סך הכל: $1 + 3 + 5 = 9$ דרגות חופש, כנדרש.

\begin{examplebox}{יישום: מומנט קוודרופול}
באינטראקציה בין קרינה לחומר, הפיתוח למולטיפולים מכיל איברים מהצורה $(\mathbf{k}\cdot\mathbf{r})^n$. האיבר הלינארי נותן את הדיפול (וקטור, $k=1$). האיבר הריבועי מערב מכפלות $x_i x_j$, ופירוק זה מאפשר להפריד בין תרומות עם סימטריות שונות (כמו המעבר הקוודרופולי החשמלי, הקשור לחלק הסימטרי חסר העקבה).
\end{examplebox}